\documentclass[../main.tex]{subfiles}
\graphicspath{{\subfix{../img/}}}
\begin{document}
	\subsubsection{Introduzione}
	Premessa: Si ricordi che il \textbf{diodo è un componente non lineare}.\\\\
	Nella scorsa sezione abbiamo analizzato il modello dinamico del diodo ed abbiamo notato che il suo comportamento varia fortemente dalla fase di spegnimento (che necessità di un tempo $t_{\text{OFF}}$ detto tempo di storage non nullo) alla fase di accensione (che richiede invece un tempo nullo).\\\\
	Abbiamo inoltre notato che, nonostante il tempo di storage sia indipendente dalla resistenza di carico, il recupero della condizione asintotica dipende proprio dalla resistenza di carico.\\\\
	Infine abbiamo mostrato che se il tempo di storage è paragonabile al periodo del segnale di ingresso, il diodo non riesce a spegnersi e quindi è come se il diodo stesso non fosse presente (il condensatore in parallelo al diodo lo cortocircuita). \\
	Allora esistono un periodo minimo ed una frequenza massima oltre i quali il modello non è più utilizzabile.
	\subsection{Diodo $npn$}
	Fino ad ora abbiamo ragionato utilizzando il diodo a giunzione $pn$ in cui la corrente circola prevalentemente in una sola direzione. Con il diodo abbiamo costruito sia gli AND sia gli OR. Resta però da realizzare il NOT che è piuttosto complesso da realizzare con i diodi.
	Ricordiamo che tutti i ragionamenti fatti sino ad ora per il diodo sono validi con alcune limitazioni:
	\begin{itemize}
		\item Lo scostamento alla regione di equilibrio deve essere ridotto
		\item Il diodo deve essere sufficientemente grande da poter essere partizionato in due regioni neutre ai margini ed in due regioni svuotate limitrofe alla giunzione
	\end{itemize}
	Costruiamo ora una struttura più complessa: la giunzione $npn$ che aggangia una regione $n$ ad ogni lato della regione $p$ (Figura 12.1).
	\begin{center}
		\begin{figure}[h]
			\centering
			\includegraphics[scale=0.35]{50.png}
			\caption{Dispositivo composto da due diodi}
			\label{fig:im-50}
		\end{figure}
	\end{center}
	I terminali sono indicati come emettitore ($E$), base ($B$) e collettore ($C$).\\\\
	Possiamo pensare di analizzare il sistema appena costruito con quanto già sappiamo del diodo a giunzione $p_n$ poiché si tratta, in effetti, di un circuito composto da due diodi.\\\\
	Detta $w$ la \textbf{distanza tra le due regioni} $n$, abbiamo però due comportamenti differenti per la situazione in cui il diodo $D2$ è spento:
	\begin{enumerate}
		\item Per $w$ superiore a qualche micron, la corrente di collettore $I_C^*$ è trascurabile (allora il circuito si riduce al solo diodo $D1$ che abbiamo già studiato in precedenza):
		$$\begin{cases} V_{BE} > 0 \\ V_{BC} < 0 \end{cases} \implies \begin{cases} I_1 > 0 \\ I_2 = 0 \end{cases} \implies \begin{cases} I_B = I_E = I_1 \\ I_C^* = 0 \end{cases}$$
		Ricordando per i calcoli:
		$$V_{BE} - V_{BC} - V_{CE} = 0$$
		$$I_B + I_C = I_E \quad \text{Kirchoff}$$
		$\begin{cases}
			BE = \text{ON} \\
			BC = \text{OFF}
		\end{cases}$
		\item Per $w$ inferiore a qualche micron, la corrente di collettore $I_C^*$ è paragonabile alla corrente di emettitore $-I_E$ e la corrente di base $I_B$ è molto inferiore alle altre due correnti:
		$$\begin{cases} V_{BE} > 0 \\ V_{BC} < 0 \end{cases} \implies \begin{cases} I_1 > 0 \\ I_2 = 0 \end{cases} \implies \begin{cases} I_C^* = -I_E \\ I_B \ll I_C^* \end{cases}$$
	\end{enumerate}
	Le Figure \ref{fig:im-51} mostrano i diversi comportamenti del dispositivo in funzione della dimensione di $w$.
	\begin{center}
		\begin{figure}[h]
			\centering
			\includegraphics[scale=0.55]{51.png}
			\caption{Diversi comportamenti nella regione svuotata. \\ Una base di dimensioni sufficientemente elevate (a sinistra) consente di studiare il dispositivo come fatto in precedenza.}
			\label{fig:im-51}
		\end{figure}
	\end{center}
	Il modello precedente nel secondo caso non funziona poiché, se $w$ è troppo piccolo, la regione $p$ non è più partizionale in regioni svuotate e regioni neutre. \\
	Quando $w$ è di dimensioni inferiori a pochi micron, le due regioni $n$ hanno un effetto complementare sulla regione svuotata: gli elettroni hanno probabilità sempre maggiori di sopravvivere al passaggio nella base.\\\\
	In questo caso, la giunzione polarizzata in inversa trasporta molta carica grazie al supporto della giunzione polarizzata in diretta che pone nella base portatori minoritari recuperati poi dal collettore.\\\\
	Questo effetto è detto effetto transistor ed è alla base dei componenti detti transistori. \\
	In particolare il componente appena realizzato è detto transistore bipolare a giunzione (o BJT, bipolar junction transistor) ed è indicato con il simbolo di Figura \ref{fig:im-52}.
	\begin{center}
		\begin{figure}[h]
			\centering
			\includegraphics[scale=0.35]{52.png}
			\caption{Simboli grafici per la rappresentazione di un transistore npn (a sinistra) o pnp (a destra) in un circuito elettrico}
			\label{fig:im-52}
		\end{figure}
	\end{center}
	\subsection{Modello di Ebers \& Moll}
	L’emettitore è così detto perché immette portatori nella base ove si trovano le lacune. \\
	Una parte sperabilmente piccola degli elettroni immessi nella base, ricombinerà con le lacune lì presenti. \\
	Tuttavia il collettore, recupererà una gran parte degli elettroni immessi nella base dall’emettitore.
	I flussi di lacune ed elettroni sono mostrati in Figura \ref{fig:im-53}.
	\begin{center}
		\begin{figure}[h]
			\centering
			\includegraphics[scale=0.35]{53.png}
			\caption{Flussi di elettroni e lacune all’interno del modello di Ebers \& Moll}
			\label{fig:im-53}
		\end{figure}
	\end{center}
	Detta $I_t$ la corrente dovuta all’effetto transistor, possiamo utilizzare il modello studiato per il diodo e completarlo con l’aggiunta del termine $I_t$, definito, dalle legge di Kirchhoff, come (Che dall'equazione qui sotto notiamo avere 2 componenti):
	$$I_t = \underbrace{I_S \left( e^{\frac{V_{BE}}{V_T}} - 1 \right)}_{\text{Componente 1}} - \underbrace{I_S \left( e^{\frac{V_{BC}}{V_T}} - 1 \right)}_{\text{Componente 2}}$$
	Parimenti, possiamo definire le correnti sui due diodi, come:
	$$I_1 = I_{BES} \left( e^{\frac{V_{BE}}{V_T}} - 1 \right) \qquad \left(I_{BES} = I_{SE} - I_S\right)$$
	$$I_2 = I_{BCS} \left( e^{\frac{V_{BC}}{V_T}} - 1 \right) \qquad \left(I_{BCS} = I_{SC} + I_B\right)$$
	Notiamo che, siccome la corrente di collettore è elevata quando si è in regione di polarizzazione diretta, assumeremo che sia $I_C = -I_C^*$.\\\\
	I grafici dei due possibili BJT ($npn$ e $pnp$), sono riportati in Figura \ref{fig:im-52}. \\
	Utilizzeremo prevalentemente il modello $npn$.\\\\
	Il modello riportato in Figura \ref{fig:im-54} è chiamato Modello di Ebers \& Moll e funziona sufficientemente bene per le nostre applicazioni.
	\begin{center}
		\begin{figure}[h]
			\centering
			\includegraphics[scale=0.35]{54.png}
			\caption{Grafico del modello di Ebers \& Moll}
			\label{fig:im-54}
		\end{figure}
	\end{center}
	Il modello è descritto da:
	$$\begin{cases} 
		I_B = I_{BES} \left( e^{\frac{V_{BE}}{V_T}} - 1 \right) + I_{BCS} \left( e^{\frac{V_{BC}}{V_T}} - 1 \right) \\ 
		I_C = I_S \left( e^{\frac{V_{BE}}{V_T}} - 1 \right) - (I_{BCS} + I_S) \left( e^{\frac{V_{BC}}{V_T}} - 1 \right) \\ 
		I_E = (I_{BES} + I_S) \left( e^{\frac{V_{BE}}{V_T}} - 1 \right) - I_S \left( e^{\frac{V_{BC}}{V_T}} - 1 \right) 
	\end{cases}$$
	Si dice che il modello di \textbf{Ebers \& Moll} è in regione normale (o attiva diretta) quando si ha:
	$$\begin{cases} V_{BE} > 0 \\ V_{BC} < 0 \end{cases} \implies \begin{cases} V_B > V_E \\ V_B < V_C \end{cases} \implies V_E < V_B < V_C$$
	Allora possiamo controllare la corrente tra collettore ed emettitore agendo sulla tensione della base. \\
	Se $V_B$ è bassa, la corrente tra collettore ed emettitore è bassa; se $V_B$ è alta, la corrente tra collettore ed emettitore è alta.
	\subsection{Semplificazioni}
	Il modello di Ebers \& Moll è \textbf{notevolmente più complesso} rispetto al modello usato per il diodo a giunzione $pn$. \\
	Possiamo tuttavia semplificarci un poco la vita sfruttando le leggi di Kirchhoff:
	$$I_C + I_B = I_E$$
	$$V_{BE} = V_{BC} + V_{CE} \iff V_B - V_E = V_B - V_C + V_C - V_E$$
		Possiamo inoltre ottenere ulteriori semplificazioni a seconda della regione di funzionamento.\\\\
	Poniamoci in regione attiva diretta (funzionamento normale del transistore). \\
	Per ipotesi, abbiamo:
	$$\begin{cases} V_{BE} > 0 \\ V_{BC} < 0 \end{cases} \implies V_{BE} > V_{BC} \implies \left( e^{\frac{V_{BE}}{V_T}} - 1 \right) \gg \left( e^{\frac{V_{BC}}{V_T}} - 1 \right)$$
	ed allora possiamo semplificare l’espressione delle correnti trascurando l’esponenziale di $V_{BC}$:
	$$\begin{cases}I_{B}\approx I_{BES}\left(e^{\frac{V_{BE}}{V_{T}}}-1\right)\\I_{C}\approx I_{S}\left(e^{\frac{V_{BE}}{V_{T}}}-1\right)\\I_{E}\approx(I_{S}+I_{BES})\left(e^{\frac{V_{BE}}{V_{T}}}-1\right)\end{cases}$$
	Allora tutte le correnti dipendono solamente dalla tensione applicata alla base. \\
	Inoltre possiamo semplicemente calcolare la corrente di collettore nota la corrente di base come (\textbf{Pedice F: Forward; Pedice R: Opposto(Reverse)}):
	$$I_C = \beta_F I_B$$
	dove $\beta_F$ è una costante definita dal rapporto:
	$$\beta_F = \frac{I_C}{I_B}$$
	Analogamente abbiamo:
	$$\alpha_F = \frac{I_C}{I_E}$$
	ed $\alpha_F < 1$ prende il nome di efficienza di emettitore poiché vale:
	$$I_C = \alpha_F I_E$$
	È inoltre evidente che $\alpha_F$ e $\beta_F$ sono legati dalla relazione:
	$$\alpha_F = \frac{\beta_F}{\beta_F + 1} \implies \lim_{\beta_F \to \infty} \frac{\beta_F}{\beta_F + 1} = 1$$
	\begin{center}
		oppure
	\end{center}
	$$\beta_F = \frac{\alpha_F}{1 - \alpha_F}$$
	Allora il modello si semplifica ulteriormente:
	$$\begin{cases}I_{B}\approx I_{BES}\left(e^{\frac{V_{BE}}{V_{T}}}-1\right)\\I_{C}\approx\beta_{F}I_{B}\\I_{E}\approx I_{C}+I_{B}\end{cases}$$
	e lo schema elettrico è riportato in Figura \ref{fig:im-55}.\\\\
	Parimenti, semplificando con $\alpha_F$ è:
	$$\begin{cases}I_B\approx I_C+I_E<<I_c\\I_C\approx\alpha_FI_E \sim I_E\\I_E\approx(I_S+I_{BES})\left(e^{\frac{V_{BE}}{V_T}}-1\right)&\end{cases}$$
	che si può rappresentare come in Figura \ref{fig:im-56}.\\\\
	Dunque, in polarizzazione diretta, il transistore ha una complessità di poco superiore a quella del diodo. \\
	Allora, siccome l’equazione esponenziale è quella del diodo, possiamo utilizzare nuovamente la semplificazione vista per il diodo ed il relativo modello a soglia.
	\begin{figure}[h]
		
		\begin{subfigure}{0.5\textwidth}
			\includegraphics[width=0.9\linewidth]{55.png} 
			\caption{Semplificazione del modello di Ebers \& Moll con l’uso del parametro $\beta_F$}
			\label{fig:im-55}
		\end{subfigure}
		\begin{subfigure}{0.5\textwidth}
			\includegraphics[width=0.9\linewidth]{56.png}
			\caption{Semplificazione del modello di Ebers \& Moll con l’uso del parametro $\alpha_F$}
			\label{fig:im-56}
		\end{subfigure}
		
		\caption{}
		\label{fig:image7}
	\end{figure}
	\subsection{Regione $V_{BE} < 0$ e $V_{BC} < 0$}
	In questa regione, dal momento che entrambi gli esponenziali sono molto minori di 1, il modello si semplifica in:
	$$\begin{cases}I_B\approx-I_{BES}-I_{BCS}\\I_C\approx-I_{BCS}\\I_E\approx-I_{BES}\end{cases}$$
	L’effetto transistore in questa regione di funzionamento non si manifesta. \\
	Il modello circuitale semplificato è riportato in Figura \ref{fig:im-57}.
	\begin{center}
		\begin{figure}[h]
			\centering
			\includegraphics[scale=0.35]{57.png}
			\caption{Semplificazione del modello di Ebers \& Moll con per la regione di interdizione}
			\label{fig:im-57}
		\end{figure}
	\end{center}
	Questa regione è detta regione di interdizione. \\
	Per semplicità spesso si dice che il transistore è spento.
	\subsection{Regione $V_{BE} < 0$ e $V_{BC} > 0$}
	Ragionando analogamente a quanto fatto in precedenza per la regione attiva diretta, notiamo che l’esponenziale associato a $V_{BE}$ è trascurabile rispetto all’esponenziale associato a $V_{BC}$.
	\subsubsection{Analisi Qualitativa}
	Ci troviamo in questo caso nella regione duale rispetto alla regione normale. \\
	Il modello si semplifica in:
	$$\begin{cases} I_B = I_{BCS} \left( e^{\frac{V_{BC}}{V_T}} - 1 \right) \\ I_C = -(I_S - I_{BCS}) \left( e^{\frac{V_{BC}}{V_T}} - 1 \right) \\ I_E = -I_S \left( e^{\frac{V_{BC}}{V_T}} - 1 \right) \end{cases}$$
	Questa regione è detta regione attiva inversa. \\
	Possiamo dire che il collettore e l’emettitore si scambiano i ruoli rispetto alla regione attiva diretta. \\
	Analogamente alla regione attiva diretta, possiamo ricavare le costanti:
	$$\frac{I_E}{I_B} = -\beta_R$$
	$$\frac{I_E}{I_C} = \alpha_R$$
	da cui ricavare le relazioni:
	$$I_E = -\beta_R I_B$$
	$$I_E = \alpha_R I_C$$
	$$\alpha_R = \frac{\beta_R}{\beta_R - 1}$$
	$$\beta_R = \frac{\alpha_R}{1 - \alpha_R}$$
	L’equivalente circuitale semplificato è riportato nelle Figure \ref{fig:im-58} e \ref{fig:im-59}.
	\begin{figure}[h]
		
		\begin{subfigure}{0.5\textwidth}
			\includegraphics[width=0.9\linewidth]{58.png} 
			\caption{Semplificazione del modello di Ebers \& Moll con per la regione attiva inversa. \\ La corrente del transistore è data da $I_t = \beta_R I_R$.}
			\label{fig:im-58}
		\end{subfigure}
		\begin{subfigure}{0.5\textwidth}
			\includegraphics[width=0.9\linewidth]{59.png}
			\caption{Semplificazione del modello di Ebers \& Moll con per la regione attiva diretta. \\ La corrente di transistore è $I_t = \alpha_R I_R$.}
			\label{fig:im-59}
		\end{subfigure}
		
		\caption{}
		\label{fig:image8}
	\end{figure}
	\subsubsection{Analisi Quantitativa}
	Dal punto di vista quantitativo, la simmettria non esiste. \\
	Utilizzando il transistore in regione attiva indiretta, i parametri $\beta_R$ e $\alpha_R$ sono solitamente più piccoli rispetto a $\beta_F$ e $\alpha_F$. \\
	Per tale motivo tenderemo a non usare la regione attiva indiretta.\\\\
	Questo differente comportamento è dovuto al fatto che in regione attiva diretta vogliamo che la corrente dovuta all’effetto transistor sia dominante rispetto alla corrente dovuta al diodo. \\
	Questo provoca due flussi: uno di elettroni che vengono iniettati in base; l’altro di lacune che vengono iniettate nell’emettitore. \\
	Dal collettore vengono accettati gli elettroni provenienti dall’emettitore e vengono ineiettate in base ben poche lacune.\\\\
	Allora perché $\beta_F$ sia grande è necessario che la corrente dovuta all’effetto transistore sia grande rispetto alla corrente di base e dunque deve essere $I_S \gg I_{BES}$. \\
	Questo risultato però può essere raggiunto solamente sbilanciando fortemente le polpolazioni dei portatori di carica: la concentrazione di atomi donatori in base deve quindi essere necessariamente molto maggiore rispetto alla concentrazione degli atomi accettati in base.
	$$N_{D,E} \gg N_{A,B}$$
	Inoltre è anche necessario che $w$ sia quanto più piccolo possibile.\\\\
	Per avere un valore elevato di $\beta_R$ ragioniamo in modo uguale ed opposto. \\
	Tuttavia non è possibile che il drogaggio della regione di base sia contemporaneamente molto minore sia di quello dell’emettitore sia di quello del collettore dal momento che, mentre l’emettitore è polarizzato in diretta ed ha una regione svuotata più piccola rispetto alla condizione di quiete, il collettore è polarizzato in inversa ed ha una regione svuotata più grande rispetto alla condizione di equilibrio.\\\\
	La regione svuotata della giunzione tra base e collettore è:
	$$w_{BC} = \frac{1}{N_{A,C}} \sqrt{\frac{2\varepsilon(\psi_{BC} - V_{BC})}{q\left(\frac{1}{N_{A,C}} - \frac{1}{N_{D,C}}\right)}}$$
	è varia significativamente a seconda della variazione della $V_{BC}$. \\
	La condizione in cui la regione neutra della base non esiste più (detta di condizione di punch through) è fastidiosa non possiamo più controllare il transistore con la tensione di base.\\\\
	Allora per fare in modo che la condizione di punch through non si presenti, dobbiamo ridurre la concentrazione di drogante presente nel collettore in modo che la regione svuotata si estenda solo dalla parte del collettore e non nella base.
	$$N_{A,B} \gg N_{D,C}$$
	Ma allora, riprendendo quanto detto per l’emettitore, è:
	$$N_{D,B} \gg N_{A,B} \gg N_{D,C} \implies N_{D,B} \gg N_{D,C}$$
	Questa differenza tecnologica necessaria al funzionamento del transistore impedisce al dispositivo di funzionare bene in regione attiva inversa. \\
	Tuttavia, questa stessa differenza, ci consente di realizzare facilmente le strutture npn necessarie alla realizzazione dei transistori (Figura \ref{fig:im-60}) poiché la precisione nei drogaggi non è critica. \\
	Grazie a ciò è molto facile controllare l’ampiezza della regione di base (tecnologicamente possiamo facilmente realizzare regioni di base di dimensioni pari frazioni del decimo di micron).
	\begin{center}
		\begin{figure}[h]
			\centering
			\includegraphics[scale=0.35]{60.png}
			\caption{Differenti drogaggi consentono di realizzare facilmente un transistore bipolare (evidenziato in giallo)}
			\label{fig:im-60}
		\end{figure}
	\end{center}
	\subsection{Regione $V_{BE} > 0$ e $V_{BC} > 0$}
	Vediamo ora l’ultima regione di funzionamento del diodo: in questa regione (detta regione di saturazione) entrambe le giunzioni sono polarizzate in diretta.\\
	Entrambi gli esponenziali del modello di Ebers \& Moll sono significativi e non si può semplificare.
	$$\begin{cases}
		I_B = I_{BES} \left( e^{\frac{V_{BE}}{V_T}} - 1 \right) + I_{BCS} \left( e^{\frac{V_{BC}}{V_T}} - 1 \right) \\
		I_C = I_S \left( e^{\frac{V_{BE}}{V_T}} - 1 \right) - (I_{BCS} + I_S) \left( e^{\frac{V_{BC}}{V_T}} - 1 \right) \\
		I_E = (I_{BES} + I_S) \left( e^{\frac{V_{BE}}{V_T}} - 1 \right) - I_S \left( e^{\frac{V_{BC}}{V_T}} - 1 \right)
	\end{cases}$$
	Il fatto che il modello non possa essere semplificato trova conferma nel fatto che le correnti sono una combinazione lineare degli esponenziali. \\
	Allora, dal momento che la regione di saturazione prevede entrambi i diodi accesi, sommando le espressioni ottenute per le regioni di funzionamento attiva diretta ed attiva indiretta, torniamo ad ottenere il modello completo.\\\\
	Una rappresentazione alternativa del modello è riportata in Figura \ref{fig:im-61}.
	\begin{center}
		\begin{figure}[h]
			\centering
			\includegraphics[scale=0.35]{61.png}
			\caption{Seconda rappresentazione grafica del modello di Ebers \& Moll che mostra come in saturazione esso possa essere dedotto come somma tra le semplificazioni per le regioni di funzionamento attiva diretta ed attiva indiretta}
			\label{fig:im-61}
		\end{figure}
	\end{center}
	\subsection{Regioni di Funzionamento}
	Il grafico di Figura \ref{fig:im-62} mostra le quattro ragioni di funzionamento di un BJT.
	\begin{center}
		\begin{figure}[h]
			\centering
			\includegraphics[scale=0.35]{62.png}
			\caption{Regioni di funzionamento di un BJT}
			\label{fig:im-62}
		\end{figure}
	\end{center}
	\subsection{Applicazione}
	Il transistore è il primo componente che studiamo a presentare tre morsetti.\\
	Si potrebbe dunque pensare che un tripolo qualsiasi sia descritto da sei grandezze indipendenti. \\
	Per le leggi di Kirchhoff però risulta immediato dedurre che le grandezze indipendenti siano solo quattro.
	$$\begin{cases} I_1 + I_2 + I_3 = 0 \\ V_{1,2} + V_{2,3} + V_{3,1} = 0 \end{cases} $$
	Allora possiamo descrivere il comportamento del generico tripolo di Figura \ref{fig:im-63} con due sole equazioni per le correnti e due equazioni per le tensioni.
	\begin{center}
		\begin{figure}[h]
			\centering
			\includegraphics[scale=0.35]{63.png}
			\caption{Un tripolo generico}
			\label{fig:im-63}
		\end{figure}
	\end{center}
	Per determinare il comportamento della rete abbiamo però bisogno di altre due condizioni: una per il morsetto di ingresso ed una per il morsetto di uscita che leghino le correnti con le tensioni.\\\\
	Allora mentre il diodo è descritto da una equazione in due incgnite, il transistore è descritto da due equazioni in quattro incognite.
	\subsection{Configurazioni Possibili}
	Consideriamo il circuito di Figura \ref{fig:im-64} dove la scatoletta centrale è un tripolo.
	\begin{center}
		\begin{figure}[h]
			\centering
			\includegraphics[scale=0.35]{64.png}
			\caption{Connessione in emettitore comune}
			\label{fig:im-64}
		\end{figure}
	\end{center}
	Dal momento che abbiamo quattro incognite il componente tripolare dovrà essere descritto da due equazioni in quattro incognite. \\
	Le altre due equazioni sono fornite dai componenti agganciati alle porte.
	In particolare, con la configurazione di Figura \ref{fig:im-64}, abbiamo che la $I_i$ è pari alla $I_B$ e la $I_u$ è pari alla $I_C$. \\
	Le tensioni sono analogamente $V_i = V_{BE}$ e $V_u = V_{CE}$. \\
	Siccome l’emettitore è il riferimento comune per le tensioni questa configurazione è detta connessione in emettitore comune ed è descritta dalle correnti:
	$$\begin{cases} I_B(V_{BE}, V_{CE}) \\ I_C(V_{BE}, V_{CE}) \end{cases}$$
	Altre configurazioni possibili sono la \textbf{connessione in base comune} (Figura \ref{fig:im-65}) e la \textbf{connessione in collettore comune} (Figura \ref{fig:im-66}).
	\begin{figure}[h]
		
		\begin{subfigure}{0.5\textwidth}
			\includegraphics[width=0.9\linewidth]{65.png} 
			\caption{Connessione in base comune}
			\label{fig:im-65}
		\end{subfigure}
		\begin{subfigure}{0.5\textwidth}
			\includegraphics[width=0.9\linewidth]{66.png}
			\caption{Connessione in collettore comune}
			\label{fig:im-66}
		\end{subfigure}
		
		\caption{}
		\label{fig:image9}
	\end{figure}
	Le tre differenti configurazioni hanno usi e caratteristiche piuttosto differenti e si adattano meglio ad alcune situazioni piuttosto che ad altre.
	\subsection{Connessione in emettitore comune}
	Riprendiamo il circuito di Figura \ref{fig:im-67}
	\begin{center}
		\begin{figure}[h]
			\centering
			\includegraphics[scale=0.35]{67.png}
			\caption{Connessione in emettitore comune}
			\label{fig:im-67}
		\end{figure}
	\end{center}
	Dal modello di Ebers \& Moll definito in precedenza
	$$\begin{cases}
		I_B = I_{BES} \left( e^{\frac{V_{BE}}{V_T}} - 1 \right) + I_{BCS} \left( e^{\frac{V_{BC}}{V_T}} - 1 \right) \\
		I_C = I_S \left( e^{\frac{V_{BE}}{V_T}} - 1 \right) - (I_{BCS} + I_S) \left( e^{\frac{V_{BC}}{V_T}} - 1 \right) \\
		I_E = (I_{BES} + I_S) \left( e^{\frac{V_{BE}}{V_T}} - 1 \right) - I_S \left( e^{\frac{V_{BC}}{V_T}} - 1 \right)
	\end{cases}$$
	cerchiamo di esplicitare le correnti del tripolo in esame.
	$$\begin{cases} I_B(V_{BE}, V_{CE}) \\ I_C(V_{BE}, V_{CE}) \end{cases}$$
	La $V_{BC}$ non è definita nel modello.\\
	 Fortunatamente, siccome la $V_{BC}$ non è indipendente dalle altre tensioni, possiamo scriverla come:
	$$V _ {B C} = V _ {B E} - V _ {C E} \Longrightarrow V _ {C E} = V _ {B E} - V _ {B C}$$
	\subsubsection{Ingresso}
	Allora possiamo vedere la corrente di ingresso come:
	$$I_B(V_{BE}, V_{CE}) = I_{BES} \left( e^{\frac{V_{BE}}{V_T}} - 1 \right) + I_{BCS} \left( e^{\frac{V_{BE} - V_{CE}}{V_T}} - 1 \right)$$
	Questa configurazione è però molto complessa poiché implica che la corrente $I_B$ dipende da una coppia di grandezze e quindi dovrebbe essere in un grafico tridimensionale. \\
	Dobbiamo allora ricavare un sistema più pratico.\\\\
	Una soluzione molto più pratica consiste nel tracciare una serie di curve bidimensionali parametriche fissando la $V_{CE}$ o la $V_{BE}$ in modo da ottenere una serie di funzioni di una sola variabile.
	$$I_B(V_{BE}, V_{CE}) \longrightarrow I_B(V_{BE}) \text{ con } V_{CE} = \text{costante}$$
	In condizione di funzionamento normale, ricaviamo che è:
	$$e^{\frac{V_{BE}}{V_T}} \gg e^{\frac{V_{BC}}{V_T}}$$
	$$I_B = I_{BES} \left( e^{\frac{V_{BE}}{V_T}} - 1 \right)$$
	ed allora l’intera famiglia di curve collassa sulla sola curva di Figura \ref{fig:im-68}.
	\begin{center}
		\begin{figure}[h]
			\centering
			\includegraphics[scale=0.35]{68.png}
			\caption{Caratteristica dell’ingresso di un BJT}
			\label{fig:im-68}
		\end{figure}
	\end{center}
	Se invece ci poniamo nella regione di saturazione abbiamo:
	$$I_B(V_{BE}, V_{CE}) = I_{BES} \left( e^{\frac{V_{BE}}{V_T}} \right) + I_{BCS} \left( \frac{e^{\frac{V_{BE}}{V_T}}}{e^{\frac{V_{CE}}{V_T}}} \right)$$
	$$I_B = e^{\frac{V_{BE}}{V_T}} \left( I_{BES} + \frac{I_{BCS}}{e^{\frac{V_{CE}}{V_T}}} \right)$$
	e quindi la famiglia di curve presenta andamenti esponenziali moltiplicati per una costante tanto più grande quanto più $V_{CE}$ è piccolo.
	\subsubsection{Uscita}
	\label{subsub:conn-em-com-out}
	Vediamo ora cosa succede per la corrente di uscita.\\\\
	Analogamente a quanto visto prima è:
	$$I_C(V_{BE}, V_{CE}) \longrightarrow I_C(V_{CE}) \text{ con } V_{BE} = \text{costante}$$
	$$I _ { C } = I _ { S } \left( e ^ { \frac { V _ { B E } } { V _ { T } } } - 1 \right) - \left( I _ { B C S } + I _ { S } \right) \left( e ^ { \frac { V _ { B E } - V _ { C E } } { V _ { T } } } - 1 \right)$$
	In regione di funzionamento normale, il termine dovuto alla $V_{CE}$ è trascurabile e la curva di uscita risulta essere costante per valori costanti di $V_{BE}$:
	$$I _ { C } = I _ { S } \left( e ^ { \frac { V _ { B E } } { V _ { T } } } - 1 \right) = \beta _ { F } I _ { B }$$
	Questa condizione ha senso solo quando si ha:
	$$V _ { B E } > V _ { C E } \Longrightarrow I _ { S } \left( e ^ { \frac { V _ { C E } } { V _ { T } } } - 1 \right) > I _ { C }$$
	In saturazione abbiamo invece la condizione opposta a quella per il funzionamento in regione normale. \\
	La corrente di uscita è allora:
	$$I _ { C } = I _ { S } \left( e ^ { \frac { V _ { B E } } { V _ { T } } } \right) - \left( I _ { B C S } + I _ { S } \right) \left( \frac { e ^ { \frac { V _ { B E } } { V _ { T } } } } { e ^ { \frac { V _ { C E } } { V _ { T } } } } \right) < \beta _ { F } I _ { B }$$
	dove abbiamo trascurato le unità.\\\\
	La Figura \ref{fig:im-69} mostra l’uscita del BJT.
	\begin{center}
		\begin{figure}[h]
			\centering
			\includegraphics[scale=0.35]{69.png}
			\caption{Caratteristica dell’uscita di un BJT}
			\label{fig:im-69}
		\end{figure}
	\end{center}
	Dall’immagine si nota subito che tutte le curve della regione di saturazione passano per uno stesso punto che non è l’origine. \\
	In particolare l’uscita si annulla quando risulta essere:
	$$I _ { S } \left( e ^ { \frac { V _ { B E } } { V _ { T } } } \right) - \left( I _ { B C S } + I _ { S } \right) \left( \frac { e ^ { \frac { V _ { B E } } { V _ { T } } } } { e ^ { \frac { V _ { C E } } { V _ { T } } } } \right) = 0 \Longrightarrow e ^ { \frac { V _ { B E } } { V _ { T } } } = \frac { I _ { S } + I _ { B C S } } { I _ { S } } = \frac { 1 } { \alpha _ { R } }$$
	che è una costante. \\
	Passando ora ai logaritmi risulta:
	$$V _ { C E } = V _ { T } \cdot \ln \left( \frac { 1 } { \alpha _ { R } } \right)$$
	Il punto di attraversamento dell’asse delle ascisse dunque non dipende dalla tensione $V_{BE}$.\\\\
	Fisicamente la spiegazione di questo comportamento sta nel fatto che, all’aumentare della $V_{CE}$ il flusso che dal collettore porta elettroni in base diventa sempre più importante e, per valori sufficientemente elevati di $V_{CE}$ può anche dominare sul flusso che dall’emettitore porta gli elettroni al collettore.\\\\
	La Tabella \ref{tab:table-6} riassume quanto visto fino ad ora:
	\begin{table}[ht]
		\centering
		\begin{tabular}{|c|c|c|}
			\hline
			\textbf{Porta} & \textbf{Regione} & \textbf{Caratteristica} \\
			\hline
			\textbf{Ingresso} & Diretta & Famiglia collassata su una sola curva. \\
			\hline
			& Saturazione & Curve esponenziali crescenti al decrescere di $V_{CE}$. \\
			\hline
			& Interdizione & Corrente nulla. \\
			\hline
			\textbf{Uscita} & Diretta & Valori costanti per valori costanti di $V_{BE}$. \\
			\hline
			& Saturazione & Curve esponenziali decrescenti al decrescere di $V_{CE}$. \\
			\hline
			& Interdizione & Corrente nulla. \\
			\hline
		\end{tabular}
		\caption{Caratteristiche delle varie regioni di funzionamento del BJT}
		\label{tab:table-6}
	\end{table}
	\subsection{Relazione tra ingresso e uscita}
	Consideriamo nuovamente un transistore ad emettitore comune come in Figura \ref{fig:im-70}.
	\begin{center}
		\begin{figure}[h]
			\centering
			\includegraphics[scale=0.35]{70.png}
			\caption{Un semplice circuito elettrico per lo studio della relazione ingresso-uscita del BJT}
			\label{fig:im-70}
		\end{figure}
	\end{center}
	Innanzitutto notiamo che il circuito in questione non può mai funzionare in regione normale. \\
	In regione normale infatti è:
	$$\begin{cases} V_{BE} > 0 \\ V_{BC} < 0 \end{cases} \implies I_B = I_{BES} \left( e^{\frac{V_{BE}}{V_T}} - 1 \right) > 0$$ 
	$$\begin{cases} I_C = \beta_F I_B > 0 \\ V_u = -RI_C < 0 \end{cases} \implies V_u < 0$$ 
	$$\begin{cases} V_{BE} < V_{CE} \\ V_{BE} > 0 \end{cases} \implies V_{CE} > 0$$
	Ma la $V_{CE}$ è esattamente la $V_u$. \\
	Allora la $V_u$ deve essere contemporaneamente maggiore e minore di zero. \\
	Dunque la conclusione è assurda e quindi la tesi è \textbf{assurda}.\\\\
	Una possibile soluzione per far funzionare il dispositivo in regione normale consiste nell’inserimento di una tensione di polarizzazione $V_{CC}$ in serie alla resistenza $R$ che sia più grande della caduta sulla resistenza stessa (Figura \ref{fig:im-71}).
	\begin{center}
		\begin{figure}[h]
			\centering
			\includegraphics[scale=0.35]{71.png}
			\caption{L’aggiunta del generatore $V_{CC}$ consente di far funzionare il BJT in regione normale}
			\label{fig:im-71}
		\end{figure}
	\end{center}
	Con questo sistema abbiamo infatti:
	$$V_u = V_{CC} - RI_C$$
	ed essendo $V_{CC} > RI_C$ la tensione di uscita è positiva. \\
	Si nota che, per far lavorare il dispositivo in regione normale dobbiamo consumare energia. \\
	Allora abbiamo un costo energetico.\\\\
	Per quanto riguarda la polarizzazione inversa risulta essere:
	$$\begin{cases} V_{BE} < 0 \\ V_{BC} > 0 \end{cases} \implies I_B = I_{BES} \left( e^{\frac{V_{BC}}{VT}} - 1 \right) > 0$$
	$$\begin{cases} I_E = -\beta_R I_B < 0 \\ I_E = \alpha_R I_C < 0 \end{cases} \implies I_C < 0 \implies V_u = V_{CC} + R |I_C| > 0$$
	$$\begin{cases} V_{CE} < 0 \\ V_u = V_{CE} \end{cases} \implies V_u < 0$$
	che porta nuovamente ad una conclusione \textbf{assurda}.\\\\
	Allora, imponendo che il dispositivo funzioni in regione normale, gli si impedisce di funzionare in regione attiva inversa. \\
	Dunque il modello ricavato fino ad ora fornisce comunque tutti gli strumenti necessari all’analisi.\\\\
	Riprendiamo di nuovo il circuito in figura \ref{fig:im-71}.\\\\
	Come dimostrato in precedenza, questo circuito può funzionare in tre sole regioni:
	\begin{enumerate}
		\item Regione normale
		\item Regione di interdizione
		\item Regione di saturazione
	\end{enumerate}
	Indipendentemente dalla regione di funzionamento, dalla configurazione del circuito possiamo osservare che è:
	$$V_i = V_{BE}$$
	$$V_u = V_{CE} = V_{CC} - RI_C$$
	\subsubsection{Regione di Interdizione}
	Iniziamo a ragionare ponendo il transistore nella regione di interdizione. \\
	Dal modello di Ebers \& Moll è:
	$$\begin{cases} V_{BE} < 0 \\ V_{CE} < 0 \end{cases} \implies V_{BC} < 0 \implies V_{CE} > V_{BE}$$
	da cui segue che la tensione di ingresso dovrà essere negativa.\\\\
	Sempre dal modello di Ebers \& Moll, usando le note semplificazioni, troviamo:
	$$I_C \approx I_{BCS} \approx 0$$
	da cui ricaviamo che la tensione di uscita è:
	$$V_u = V_{CC} - RI_C \implies V_u = V_{CC} > 0$$
	La condizione per cui si presenta questa situazione è $V_u > V_i$. \\
	Siccome però $V_u = V_{CC} > 0$ e $V_i < 0$, la condizione è sempre verificata.
	\subsubsection{Regione Normale}
	In regione normale è:
	$$\begin{cases} V_{BE} > 0 \\ V_{CE} < 0 \end{cases} \implies \begin{cases} V_i > 0 \\ V_u > V_i \end{cases}$$
	Applichiamo le note semplificazioni del modello di Ebers \& Moll ed otteniamo:
	$$\begin{cases} I_B = I_{BES} \left( e^{\frac{V_{BE}}{V_T}} - 1 \right) \\ I_C = I_S \left( e^{\frac{V_{BE}}{V_T}} - 1 \right) \end{cases}$$
	Allora troviamo che la tensione di uscita è:
	$$V_u = V_{CC} - RI_S I_S \left( e^{\frac{V_i}{V_T}} - 1 \right)$$
	dove abbiamo usato l’uguaglianza $V_i = V_{BE}$.
	\subsubsection{Regione di Saturazione}
	In questo caso abbiamo:
	$$\begin{cases} V_{BE} > 0 \\ V_{CE} > 0 \end{cases} \implies \begin{cases} V_i > 0 \\ V_u < V_i \end{cases}$$ 
	$$I_C = I_S \left( e^{\frac{V_{BE}}{V_T}} - 1 \right) - (I_{BCS} + I_S) \left( e^{\frac{V_{BC}}{V_T}} - 1 \right)$$ 
	$$I_C = I_S \left( e^{\frac{V_i}{V_T}} - 1 \right) - (I_S + I_{BCS}) \left( e^{\frac{V_i - V_u}{V_T}} - 1 \right)$$
	Allora la tensione di uscita è:
	$$V_u = V_{CC} - R \left[ I_S \left( e^{\frac{V_i}{V_T}} - 1 \right) - (I_S + I_{BCS}) \left( e^{\frac{V_i - V_u}{V_T}} - 1 \right) \right]$$
	che è abbastanza complessa.\\\\
	Fortunatamente possiamo risolvere l’uguaglianza sfruttando il grafico della caratteristica di uscita del BJT vista nel Paragrafo \ref{subsub:conn-em-com-out}.\\\\
	Dal momento che è:
	$$\begin{cases} V_{CE} = V_u \\ V_u = V_{CC} - RI_C \end{cases} \implies I_C = \frac{V_{CC} - V_u}{R}$$
	abbiamo l’espressione di una retta. \\
	Tale retta (solitamente detta retta di carico della rete) rappresenta il luogo dei punti $(V_u, I_C)$ che soddisfa la legge di Kirchhoff alla maglia. \\
	Ogni curva rappresenta invece il luogo dei punti che soddisfano le equazioni del modello di Ebers \& Moll per una data tensione di ingresso.\\\\
	Dal grafico di Figura \ref{fig:im-72} notiamo che l’uscita cala sempre più lentamente al crescere dell’ingresso per raggiungere asintoticamente il valore:
	$$V_T \cdot \ln \left( \frac{1}{\alpha_R} \right)$$
	\begin{center}
		\begin{figure}[h]
			\centering
			\includegraphics[scale=0.45]{72.png}
			\caption{Caratteristica di uscita di un BJT. \\ In blu è evidenziata la retta di carico.}
			\label{fig:im-72}
		\end{figure}
	\end{center}
	\subsubsection{Riassunto}
	Il grafico di Figura \ref{fig:im-73} mostra quanto ottenuto dallo studio appena effettuato.
	\begin{center}
		\begin{figure}[h]
			\centering
			\includegraphics[scale=0.45]{73.png}
			\caption{Caratteristica di ingresso-uscita di un BJT ad emettitore comune. \\ In verde la caratteristica in interdizione, in rosso la caratteristica in regione normale ed in azzurro la caratteristica in saturazione.}
			\label{fig:im-73}
		\end{figure}
	\end{center}
	Questo circuito dispone di tre regioni con comportamenti qualitativamente molto differenti: due tratti a valori approssimativamente costanti (regione di interdizione e regione di saturazione) ed un tratto di raccordo (regione normale) che può avere una pendenza anche molto elevata.\\\\
	La pendenza è data dalla derivata dell’uscita. \\
	Siccome la parte con pendenza non nulla è dovuta alla regione normale, riprendiamo l’uscita in tale regione:
	$$V_u = V_{CC} - RI_S I_S \left( e^{\frac{V_i}{V_T}} - 1 \right)$$
	e ne calcoliamo la derivata nel punto $V_i = V_{i_0}$. \\
	Risulta:
	$$DV_u = -\frac{RI_S \cdot e^{\frac{V_{i_0}}{V_T}}}{V_T} = -\frac{RI_{C_0}}{V_T}$$
	Poiché siamo in regione normale, la $V_{i_0}$ è positiva ed allora possiamo trascurare il termine unitario. \\
	Allora la corrente di collettore è:
	$$I_C = I_S \cdot e^{\frac{V_{i_0}}{V_T}}$$
	Notiamo allora che la derivata ha un valore a denominatore molto basso ed un valore a numeratore molto più elevato. \\
	Allora la pendenza è molto elevata.\\\\
	Un circuito di questo tipo si presta bene alla realizzazione di una \textbf{porta NOT}. \\
	Consideriamo ad esempio i segnali:
	$$\begin{cases} V_L = 0 \\ V_H \approx V_{CC} \end{cases}$$
	La Tabella \ref{tab:table-7} mostra quanto appena detto:
	\begin{table}[ht]
		\centering
		\begin{tabular}{|c|c|}
			\hline
			$V_i$ & $V_u$ \\
			\hline 
			\hline
			$V_L$ & $V_H$ \\
			\hline
			$V_H$ & $V_L$ \\
			\hline
		\end{tabular}
		\caption{Funzionamento dell’invertitore logico realizzato con un BJT}
		\label{tab:table-7}
	\end{table}
	Si faccia attenzione al fatto che, pur non usando la regione normale di funzionamento del BJT, il componente si comporta come un invertitore proprio per la presenza della regione normale ad elevata pendenza.\\\\
	Per lavorare in regione normale è necessario progettare opportunamente il segnale di ingresso. \\
	Infatti se a $V_i = V_{i_0}$ corrisponde $V_u = V_{u_0}$ dovrà anche essere:
	$$V_i = V_{i_0} + \Delta V_i \longleftrightarrow V_u = V_{u_0} + \Delta V_u$$
	Per $\Delta V_i \rightarrow 0$ avremo allora:
	$$\frac{\Delta V_u}{\Delta V_i} = \frac{dV_u}{dV_i} = A_V \implies dV_u = A_V dV_i$$
	dove $A_V$ detto \textbf{guadagno di tensione}. \\
	Il guadagno di tensione può assumere valori anche molto elevati. \\
	Allora, quando si presenta un segnale variabile in un intervallo ristretto di $V_i$, il segnale di uscita varia in modo analogo all’ingresso amplificato di una costante $-A_V$. \\
	Allora il dispositivo studiato si comporta come un amplificatore analogico invertente.\\\\
	Il rispetto delle condizioni di piccolo spostamento e di alta pendenza è estremamente importante. \\
	Usare segnali con uno spostamento elevato provoca una distorsione del segnale di uscita rispetto al segnale di ingresso. \\
	La scelta della pendenza inficia invece il guadagno: più elevata è la pendenza, maggiore è l’amplificazione in ampiezza.\\\\
	Un’altra caratteristica del dispositivo studiato è che, mentre il segnale di ingresso è un segnale debole in potenza ed in corrente, il segnale di uscita è un segnale più forte sia in ampiezza sia in potenza. \\
	Questo fatto è dovuto al generatore di tensione $V_{CC}$.\\\\
	Ancora, pur avendo una $I_B$ libera di crescere indefinitamente, notiamo che la $I_C$ non può avere lo stesso comportamento. \\
	Quando si entra in regione di saturazione, la $I_C$ ha raggiunto il suo valore massimo (si dice quindi che è satura). \\
	Il problema è però dato dalla corrente $I_B$. \\
	La sua crescita esponenziale porta invece ad avere correnti di ingresso totalmente irragionevoli. \\
	Una possibile soluzione consiste nell’introdurre una resistenza sulla base (Figura \ref{fig:im-74}).
	\begin{center}
		\begin{figure}[h]
			\centering
			\includegraphics[scale=0.45]{74.png}
			\caption{L’inserimento della resistenza in base consente di limitare la corrente di ingresso a valori ragionevoli}
			\label{fig:im-74}
		\end{figure}
	\end{center}
	\subsection{Modello a Soglia}
	Cerchiamo ora di ricavare un modello a soglia valido per il BJT e limitato alle tre regioni usate fino ad ora.\\\\
	In precedenza abbiamo ricavato le due caratteristiche per l’ingresso e per l’uscita.
	\subsubsection{Ingresso}
	Per l’ingresso approssimiamo tutte le curve esponenziali con una retta verticale.\\\\
	In regione di funzionamento normale ed in regione di saturazione il modello a soglia risulta essere:
	$$\begin{cases} V_{BE} = V_\gamma \\ I_B > 0 \end{cases}$$
	mentre in interdizione è:
	$$\begin{cases} V_{BE} = 0 \\ I_B < V_\gamma \end{cases}$$
	\begin{center}
		\begin{figure}[h]
			\centering
			\includegraphics[scale=0.45]{75.png}
			\caption{Modello a soglia dell’ingresso di un BJT}
			\label{fig:im-75}
		\end{figure}
	\end{center}
	\subsubsection{Uscita}
	Anche in questo caso approssimiamo gli esponenziali a rette verticali. \\
	Per la regione di funzionamento normale è:
	$$\begin{cases} V_{CE} > V_{CESat} \\ I_C = \beta_F I_B \end{cases}$$
	mentre in saturazione abbiamo:
	$$\begin{cases} V_{CE} = V_{CESat} \\ I_C < \beta_F I_B \end{cases}$$
	Solitamente assumiamo che sia: $V_{CESat} = 0,2\ \text{V}$. \\
	La $V_{CESat}$ è la differenza tra la tensione di soglia dell’emettitore e la tensione di soglia del collettore.
	\begin{center}
		\begin{figure}[h]
			\centering
			\includegraphics[scale=0.45]{76.png}
			\caption{Modello a soglia dell’uscita di un BJT}
			\label{fig:im-76}
		\end{figure}
	\end{center}
	\subsubsection{Riassunto}
	Unendo tutti i risultati del modello a soglia ricaviamo un modello complessivo:
	$$\text{Regione di interdizione} = \begin{cases} V_{BE} < V_{\gamma} \\ V_{CE} > V_{CESat} & \quad V_{BC} < V^\prime_{\gamma} \\ I_B = 0 \\ I_C = 0 \end{cases}$$
	$$\text{Regione normale} = \begin{cases} V_{BE} = V\gamma \\ V_{CE} > V_{CESat} \\ I_B > 0 \\ I_C = \beta_F I_B \end{cases}$$
	$$\text{Regione di saturazione} = \begin{cases} V_{BE} = V_{\gamma} \\ V_{CE} = V_{CESat} \\ I_B = 0 \\ I_C < \beta_F I_B \end{cases}$$
	\subsection{Invertitore RTL}
	\label{sub:inv-rtl}
	Consideriamo ora il circuito di Figura \ref{fig:im-77} e cerchiamo la caratteristica ingresso-uscita sfruttando il modello a soglia.
	\begin{center}
		\begin{figure}[h]
			\centering
			\includegraphics[scale=0.45]{77.png}
			\caption{Un invertitore RTL.}
			\label{fig:im-77}
		\end{figure}
	\end{center}
	\subsubsection{Considerazioni iniziali}
	Indipendentemente dalla regione di funzionamento possiamo dire:
	$$\begin{cases} I_B = \frac{V_i - V_{BE}}{R_B} \\ I_C = \frac{V_{CC} - V_u}{R_C} \\ V_{CE} = V_u \end{cases}$$
	\subsubsection{Regione di Interdizione}
	Dal modello a soglia sappiamo che è:
	$$\begin{cases} I_B = 0 \\ I_C = 0 \end{cases} \implies \begin{cases} V_i = V_{BE} \\ V_u = V_{CC} \end{cases}$$
	ma siccome è $V_i = V_{BE}$ e $V_u = V_{CE}$ risulta anche:
	$$\begin{cases} V_i < V_\gamma \\ V_{CC} = V_{CE} > V_{CESat} \end{cases}$$
	\subsubsection{Regione Normale}
	In regione normale è:
	$$\begin{cases} I_C = \beta_F I_B \\ I_B = \frac{V_i - V_\gamma}{R_B} \\ V_u = V_{CC} - R_C I_C \end{cases} \implies \begin{cases} V_u = V_{CC} - \frac{\beta_F R_C}{R_B} (V_i - V_\gamma) \\ V_i > V_\gamma \\ V_{CE} = V_u > V_{CESat} \end{cases}$$
	Notiamo che l’uscita è proporzionale all’ingresso secondo la funzione:
	$$V_u = V_{CC} - \frac{\beta_F R_C}{R_B} (V_i - V_\gamma)$$
	la cui derivata è:
	$$A_V = -\frac{\beta_F R_C}{R_B}$$
	che porta ad un guadagno lineare dato dal rapporto tra la resistenza di base e la resistenza di collettore.\\\\
	Per ricavare il valore di $V_i^*$, ossia del punto in cui la retta interseca la retta costante di $V_{CESat}$, imponiamo:
	$$V_i^* = V_\gamma + \frac{V_{CC} - V_{CESat}}{|A_V|}$$
	\subsubsection{regione di Saturazione}
	In questa regione ovviamente l’uscita è:
	$$V_u = V_{CESat}$$
	e vale nella regione in cui è:
	$$V_i > V_i^*$$
	La verifica è immediata dal modello a soglia. \\
	Infatti deve essere:
	$$I_B = \frac{V_i - V_\gamma}{R_B} \implies V_i > V_\gamma$$
	$$I_C = \frac{V_{CC} - V_{CESat}}{R_C} < \beta_F \cdot \frac{V_i - V_\gamma}{R_B} \implies V_i > V_\gamma + \frac{V_{CC} - V_{CESat}}{\frac{\beta_F R_C}{R_B}}$$
	Siccome è:
	$$V_i^* = V_\gamma + \frac{V_{CC} - V_{CESat}}{|A_V|}$$
	allora risulta che è:
	$$V_i > V_i^*$$
	\subsubsection{Considerazioni finali}
	Il grafico dei risultati ottenuti fino ad ora è riportato in Figura \ref{fig:im-78}.
	\begin{center}
		\begin{figure}[h]
			\centering
			\includegraphics[scale=0.45]{78.png}
			\caption{Caratteristica ingresso-uscita dell’invertitore RTL}
			\label{fig:im-78}
		\end{figure}
	\end{center}
	Il ramo di caratteristica nella regione di interdizione è ovviamente identico a quanto visto senza la resistenza $R_B$. \\
	Tale resistenza è infatti inserita per limitare la corrente ma a transistore spento la corrente è nulla e dunque la resistenza $R_B$ non interviene.\\\\
	Nella regione normale l’uscita è rappresentata da una retta. \\
	Se l’approssimazione del modello fosse buona, otterremmo un grande vantaggio da questa caratteristica poiché consentirebbe di gestire facilmente l’amplificazione dal momento che una amplificazione lineare non distorce il segnale ricevuto.\\\\
	Dalle simulazione svolte al calcolatore sul dispositivo reale possiamo notare che l’andamento ottenuto dal modello a soglia è qualitativamente identico all’andamento della caratteristica ingresso-uscita del dispositivo reale.\\\\
	Il circuito in Figura \ref{fig:im-77} è detto Invertitore RTL (\textbf{Resistor Transistor Logic}, logica a transistore e resistore) ed è un circuito fuzionante.
	\subsection{Invertitore ad accoppiamento diretto}
	Fino a questo momento il modello a soglia ha sempre funzionato perfettamente.\\\\
	Consideriamo ora il circuito di Figura \ref{fig:im-79} e cerchiamone la caratteristica ingresso-uscita sfruttando il modello a soglia.
	\begin{center}
		\begin{figure}[h]
			\centering
			\includegraphics[scale=0.45]{79.png}
			\caption{Un invertitore ad accoppiamento diretto}
			\label{fig:im-79}
		\end{figure}
	\end{center}
	\subsubsection{Considerazioni iniziali}
	Indipendentemente dalla regione di funzionamento possiamo dire:
	$$\begin{cases} I_C = \frac{V_u - V_{CC}}{R_C} \\ V_i = V_{BE} \\ V_u = V_{CE} \end{cases}$$
	\subsubsection{Regione di interdizione}
	Dal modello a soglia, in interdizione, abbiamo:
	$$\begin{cases} I_B = 0 \\ I_C = 0 \end{cases} \implies \begin{cases} V_i = V_{BE} \\ V_u = V_{CC} \end{cases}$$
	ma siccome è $V_i = V_{BE}$ e $V_u = V_{CE}$ risulta anche:
	$$\begin{cases} V_i < V_\gamma \\ V_{CC} = V_{CE} > V_{CESat} \end{cases}$$
	\subsubsection{Regione normale}
	In regione normale è:
	$$V_{BE} = V_i = V_\gamma$$
	Ma perché sia vero deve anche essere:
	$$I_C = \beta_F I_B > 0 \implies \frac{V_u - V_{CC}}{R_C} > 0 \implies V_u > V_{CC}$$
	e, contemporaneamente, deve valere:
	$$V_{CE} = V_u > V_{CESat}$$
	da cui è:
	$$V_{CESat} < V_u < V_{CC}$$
	che rappresenta una retta veritcale.
	\subsubsection{Regione di saturazione}
	La regione di saturazione è individuata dalle condizioni:
	$$\begin{cases} V_u = V_{CE} = V_{CESat} \\ V_i = V_{BE} = V_\gamma \end{cases}$$
	che sono soddisfatte in un unico punto.
	\subsubsection{Considerazioni finali}
	Il grafico dei risultati ottenuti fino ad ora è riportato in Figura \ref{fig:im-80}.
	\begin{center}
		\begin{figure}[h]
			\centering
			\includegraphics[scale=0.45]{80.png}
			\caption{Caratteristica ingresso-uscita dell’invertitore ad accoppiamento diretto}
			\label{fig:im-80}
		\end{figure}
	\end{center}
	Il ramo di caratteristica nella regione di interdizione è ovviamente identico a quanto visto con la resistenza $R_B$. \\
	Tale resistenza è infatti inserita per limitare la corrente ma a transistore spento la corrente è nulla e dunque la resistenza $R_B$ non interviene.\\\\
	La retta verticale della regione normale presenta un guadagno infinito. \\
	Il guadagno infinito obbliga però a trattare un segnale di ingresso di ampiezza infinitesima provocando così la perdita dell’informazione.\\\\
	Il peggio tuttavia è rappresentato dal fatto che l’intera regione di saturazione implode in un unico punto. \\
	In questo caso allora il modello a soglia non fornisce una previsione sufficientemente accurata. \\
	La caratteristica ingresso-uscita nella regione di saturazione infatti è completamente sconosciuta.
	\subsection{Usabilità del modello a soglia}
	La differenza tra i circuiti di Figura \ref{fig:im-77} e \ref{fig:im-79} è molto limitata: la resistenza $R_B$ sulla maglia di ingresso presente nel primo e non nel secondo. \\
	Siccome la maglia di uscita è identica nei due circuiti dovremo concentrare la nostra analisi sulla maglia di ingresso. \\
	Inoltre, dal momento che il modello a soglia ha fallito, usiamo il modello di Ebers \& Moll.\\\\
	Consideriamo il circuito di Figura \ref{fig:im-77}.
	La corrente di base è:
	$$I_B = I_{BES} \left( e^{\frac{V_{BE}}{V_T}} - 1 \right) \implies V_{BE} = V_T \cdot \ln \left( \frac{I_B}{I_{BES}} + 1 \right)$$
	ma può anche essere espressa come:
	$$I_B = \frac{V_i - V_{BE}}{R_B} \implies V_i = V_{BE} + \underbrace{R_B I_B}_{V_R}$$
	Dunque la tensione di ingresso è una somma di due termini, entrambi dipendenti dalla corrente $I_B$. \\
	Allora possiamo pensare di variare la corrente $V_i$ da un valore $V_{i0}$ ad un valore $V_{i0} + \Delta V_i$. \\
	Conseguentemente avremo una variazione di $I_B$ da $I_{B0}$ ad $I_{B0} + \Delta I_B$.
	$$V_i : V_{i0} \rightarrow V_{i0} + \Delta V_i$$
	$$I_B : I_{B0} \rightarrow I_{B0} + \Delta I_B$$
	Allora una variazione della tensione di ingresso può essere vista come:
	$$\Delta V_i = \Delta V_R + \Delta V_{BE}$$
	Siccome il termine dato da $\Delta V_{BE}$ è molto minore rispetto al termine $\Delta V_R$, può essere trascurato:
	$$\Delta V_i \approx \Delta V_R$$
	Trascurare la caduta di tensione sulla giunzione è trascurabile solo poiché la caduta di tensione sulla resistenza è presente. \\
	Se la resistenza $R_B$ non fosse presente, la caduta dovuta alla giunzione non sarebbe più trascurabile. \\
	Questa situazione è riportata in Figura \ref{fig:im-81}.
	\begin{center}
		\begin{figure}[h]
			\centering
			\includegraphics[scale=0.45]{81.png}
			\caption{La componente esponenziale dovuta alla giunzione (in verde) è trascurabile solo in presenza di una resistenza (in azzurro). \\In giallo la variazione $\Delta V_{BE}$; in grigio la variazione $\Delta V_R$.}
			\label{fig:im-81}
		\end{figure}
	\end{center}
	Allora possiamo utilizzare il modello a soglia quando abbiamo circuiti che presentano, in serie alla giunzione, una o più resistenze. \\
	Le resistenze infatti presentano cadute di tipo ohmico che rendono trascurabili le cadute dovute alle giunzioni.
	\subsection{Qualità dell'invertitore RTL}
	Nella sezione \ref{sub:inv-rtl} abbiamo dimostrato che il circuito di Figura \ref{fig:im-77} rappresenta un invertitore. \\
	Vogliamo ora studiarne la qualità.
	\subsubsection{Invertitore}
	Prima di tutto è necessario accertarci che il circuito sia effettivamente un invertitore logico. \\
	Infatti, scelti due valori:
	$$\begin{cases} V_L = V_{CESat} \\ V_H = V_{CC} \end{cases}$$
	mentre per il valore $V_{CESat}$ siamo certi che esso sia minore di $V_\gamma$, per il valore alto $V_{CC}$ non sappiamo se valga:
	$$V_{CC} > V_i^*$$
	Questa condizione è vera solo quando abbiamo:
	$$|A_V| > \frac{V_{CC} - V_{CESat}}{V_{CC} - V_\gamma}$$
	che, fortunatamente, non è un vincolo critico. \\
	Questo risultato, però, mostra che, anche per avere un invertitore, dobbiamo avere un guadagno elevato.
	\subsubsection{Resistenza ai disturbi (Effeto del rumore)}
	\label{subsub:res-dist}
	Una volta assicuratici di disporre di un invertitore logico, dobbiamo controllarne la resistenza ai disturbi in modo da poterlo effettivamente utilizzare in una rete logica o come amplificatore.\\\\
	Possiamo interpretare il rumore come una piccola variazione rispetto al segnale che si desidera rappresentare. \\
	L’effetto di questa variazione è estremamente diverso in base alla regione in cui si lavora.
	\paragraph{Segnale analogico}
	Il segnale analogico sfrutta il tratto ad elevata pendenza della caratteristica.
	$$|A_V| = \frac{dV_u}{dV_i} = \lim_{\Delta V_i \rightarrow 0} \frac{\Delta V_u}{\Delta V_i} > 1$$
	Inoltre, per quanto riguarda i segnali analogici, una delle ipotesi fondamentali formulate all’inizio dell’analisi del BJT, era quella di piccolo scostamento.\\\\
	In presenza di rumore allora l’amplificatore amplifica sia il segnale analogico ricevuto sia il rumore rendendo di fatto indistinguibile la variazione dell’uscita voluta (cioè quella associata al segnale) da quella non voluta (cioè quella associata al rumore).
	\paragraph{Segnale digitale}
	Il segnale digitale sfrutta i tratti a pendenza quasi nulla.
	$$|A_V| = \frac{dV_u}{dV_i} = \lim_{\Delta V_i \to 0} \frac{\Delta V_u}{\Delta V_i} = 0$$
	In questo caso, un rumore che scosti di poco il segnale di ingresso dal suo valore nominale non provoca variazioni significative sul valore di uscita.\\\\
	La rete digitale è allora più immune ai disturbi rispetto all’amplificatore poiché è in grado di distinguere il segnale dal rumore. \\
	Il circuito di Figura \ref{fig:im-77} allora funziona meglio come circuito digitale rispetto a quanto può fare come amplificatore.
	\subsubsection{Immunità ai disturbi}
	\paragraph{Valori fondamentali}
	Abbiamo scoperto che l’invertitore RTL è in grado, quando usato come porta logica binaria, di discriminare tra il segnale ed il rumore. \\
	Questa discriminazione è fattibile solo quando il rumore (indicato con $N$ per noise) non supera una certa ampiezza. \\
	Cioè vogliamo individuare il massimo valore di $N$ tale per cui sia
	$$V_i^{**} = V_L + N \longrightarrow V_u = V_H$$
	Dalla caratteristica possiamo facilmente individuare il valore massimo del rumore come:
	$$N_{Max} = V_i^{**} - V_L$$
	Ancora dalla caratteristica possiamo identificare facilmente tale valore limite come $V_\gamma$. In generale indicheremo tale punto con:
	$$V_{iL \ Max}$$
	e quindi possiamo ricavare:
	$$N_{Max} = V_{iL \ Max} - V_L$$
	che nel nostro caso specifico, con valori tipici di, è:
	$$N_{Max} = V_\gamma - V_{CESat} = 0,75 - 0,2 = 0,55\ \text{V}$$
	Analogamente possiamo individuare un valore di rumore tale da non modificare l’uscita quando in ingresso abbiamo $V_H$.
	$$V_i^{***} = V_H + N \longrightarrow V_u = V_L$$
	Di nuovo, dalla caratteristica, ricaviamo che il massimo rumore sovrapponibile a $V_H$ è:
	$$N_{Max} = V_{iH \ Min} = V_\gamma + \frac{V_{CC} - V_{CESat}}{|A_V|}$$
	$$N_{Max} = V_H - V_i^* = V_{CC} - V_\gamma - \frac{V_{CC} - V_{CESat}}{|A_V|} = 3,77\ \text{V}$$
	Si nota che questo margine dipende dal guadagno. \\
	Maggiore è il guadagno, maggiore è l’immunità ai disturbi per il valore alto.\\\\
	Possiamo inoltre identificare l’escursione (\textbf{Logic Swing}) come:
	$$\Delta V_u = V_H - V_L$$
	ed il massimo errore tollerabile rispetto al valore basso come:
	$$N_{ML} = V_{iL \ Max} - V_L$$
	mentre per il valore alto è:
	$$N_{MH} = V_H - V_{iH \ Min}$$
	Con i valori tipici notiamo che $N_{ML}$ è significativamente più piccolo rispetto a $N_{MH}$. \\
	In fase di progetto dovremo allora considerare $N_{ML}$ come parametro di resistenza. In generale sarà:
	$$N_M = \min \left\{N_{ML}, N_{MH}\right\}$$
	Imponendo che entrambi i margini siano positivi e facendo la somma delle due relazioni otteniamo:
	$$N_{MH} + N_{ML} = V_H - V_{iH \ Min} + V_{iL \ Max} - V_L > 0$$
	da cui possiamo ricavare:
	$$\underbrace{V_H - V_L}_{\Delta V_u} > \underbrace{V_{iH \ Min} - V_{iL \ Max}}_{\Delta V_i}$$
	Ma $V_H - V_L$ è l’escursione mentre $V_{iH \ Min} - V_{iL \ Max}$ è l’escursione dell’ingresso. \\
	Allora:
	$$\frac{\Delta V_u}{\Delta V_i} > 1 \implies |A_V| > 1$$
	e quindi, per avere un buon invertitore, dobbiamo necessariamente realizzare un buon amplificatore. \\
	La Figura \ref{fig:im-82} mostra graficamente i margini appena individuati.\\\\
	Siccome la somma tra $N_{ML}$ e $N_{MH}$ è limitata, se aumentiamo uno, l’atro deve diminuire. \\
	La condizione ottima di progetto (Figura \ref{fig:im-83}) è:
	$$N_{ML} = N_{MH}$$
	\begin{figure}[h]
		
		\begin{subfigure}{0.5\textwidth}
			\includegraphics[width=0.9\linewidth]{82.png} 
			\caption{Rappresentazione grafica dei margini di immunità ai disturbi}
			\label{fig:im-82}
		\end{subfigure}
		\begin{subfigure}{0.5\textwidth}
			\includegraphics[width=0.9\linewidth]{83.png}
			\caption{La condizione ottima di progetto consiste nell’individuazione del miglior equilibrio possibile tra $N_{ML}$ e $N_{MH}$}
			\label{fig:im-83}
		\end{subfigure}
		
		\caption{}
		\label{fig:image10}
	\end{figure}
	\paragraph{Esempio Pratico}
	Consideriamo il circuito di Figura \ref{fig:im-84} e cerchiamo di individuarne l’immunità ai disturbi.
	\begin{center}
		\begin{figure}[h]
			\centering
			\includegraphics[scale=0.45]{84.png}
			\caption{L’aggiunta del diodo aumenta l’immunità ai disturbi della rete}
			\label{fig:im-84}
		\end{figure}
	\end{center}
	Dal momento che in serie alle giunzioni abbiamo un caricoresistivo possiamo utilizzare i modelli a soglia. \\
	Per ridurre le regioni di studio, osserviamo preliminarmente che:
	$$I_D = I_B$$
	Allora se il diodo è acceso deve valere:
	$$I_D > 0 \Longrightarrow I_B > 0$$
	e quindi il transistore deve essere acceso. Parimenti, se il transistore è acceso, il diodo non può essere spento. \\
	Allora dobbiamo valutare solo tre casi contro i sei indicati dal modello a soglia.
	\subparagraph{Regione di interdizione}
	In questa regione è:
	$$\begin{cases} I_C = 0 \\ V_u = V_{CC} - R_C I_C \end{cases} \quad \Longrightarrow V_u = V_{CC} $$
	$$ \begin{cases} V_{BE} < V_\gamma \\ V_D < V_\gamma \end{cases}$$
	$$\begin{cases} V_i - R_B I_B - V_D - V_{BE} = 0 \\ I_B = 0 \end{cases} \implies V_i < 2V_\gamma$$
	\subparagraph{Regione normale} 
	In regione normale abbiamo:
	$$\begin{cases} V_{BE} > V_\gamma \\ V_D > V_\gamma \end{cases}$$
	$$\begin{cases} I_B = \frac{V_i - 2V_\gamma}{R_B} \\ I_C = \beta_F I_B \\ V_u = V_{CC} - R_C I_C \end{cases} \implies V_u = V_{CC} - \frac{\beta_F R_C}{R_B} (V_i - 2V_\gamma)$$
	La $V_u$ è una retta della stessa pendenza della retta del circuito di Figura \ref{fig:im-77}. \\
	Tale tratto si estende in:
	$$2V_\gamma < V_i < 2V_\gamma + \frac{V_{CC} - V_{CESat}}{|A_V|}$$
	\subparagraph{Regione di saturazione} 
	In questo caso, ovviamente è:
	$$\begin{cases} V_u = V_{CESat} \\ V_i > V_i^* \end{cases}$$
	dove:
	$$V_i^* = 2V_\gamma + \frac{V_{CC} - V_{CESat}}{|A_V|}$$
	\subparagraph{Considerazioni finali} 
	La caratteristica è perfettamente identica alla caratteristica del circuito di Figura \ref{fig:im-77} ma è stata traslata in avanti di $V_\gamma$. \\
	Allora i margini si sono spostati. \\
	Risultano infatti:
	$$N_{ML} = 2V_\gamma - V_{CESat} = 1,3 \text{ V}$$
	$$N_{MH} = V_{CC} - 2V_\gamma - \frac{V_{CC} - V_{CESat}}{|A_V|} = 2,02 \text{ V}$$
	$$N_M = 1,3 \text{ V}$$
	L’aggiunta del diodo ha quindi più che raddoppiato il margine di resistenza al rumore. \\
	In linea di principio possiamo pensare di aggiungere diodi per spostare la caratteristica. \\
	Queste soluzioni sono praticate ma senza portarle a conseguenze estreme: l’aggiunta di troppi diodi infatti provoca un peggioramento delle prestazioni dinamiche.
	\subsection{Generalizzazioni}
	Quanto visto nella sezione \ref{subsub:res-dist} è estremamente limitato alla sola porta logica NOT. \\
	Dobbiamo ora trovare un modo per generalizzare quanto detto a più porte logiche (come le porte AND e OR).\\\\
	Dobbiamo inoltre generalizzare i margini di rumore a caratteristiche che non si basino sulla caratteristica specifica dell’invertitore RTL.
	\subsection{Generalizzazioni a funzioni a più ingressi}
	Consideriamo il circuito di Figura \ref{fig:im-85}
	\begin{center}
		\begin{figure}[h]
			\centering
			\includegraphics[scale=0.45]{85.png}
			\caption{Circuito per la realizzazione di una porta logica NOR}
			\label{fig:im-85}
		\end{figure}
	\end{center}
	Per semplificare l’analisi supponiamo che sia:
	$$V_1, V_2, V_u \in \left\{V_L, V_H \right\}$$
	Allora possiamo riferirci alla tabella di verità riportata in Tabella \ref{tab:table-8}.
	\begin{table}[ht]
		\centering
		\begin{tabular}{|c|c|c|}
			\hline
			$V_1$ & $V_2$ & $V_u$ \\
			\hline
			$V_L$ & $V_L$ & $V_H$ \\
			\hline
			$V_L$ & $V_H$ & $V_L$ \\
			\hline
			$V_H$ & $V_L$ & $V_L$ \\
			\hline
			$V_H$ & $V_H$ & $V_L$ \\
			\hline
		\end{tabular}
		\caption{Tabella di funzionamento del circuito di Figura \ref{fig:im-85}}
		\label{tab:table-8}
	\end{table}
	\subsubsection{$V_1 = V_L$ ; $V_2 = V_L$}
	In questo caso il BJT $T1$ è spento. \\
	Infatti, se fosse accesso, sarebbe:
	$$\begin{cases} V_{BE1} = V_\gamma \\ I_{B1} > 0 \\ I_{B1} = \frac{V_1 - V_{BE1}}{R_B} \end{cases} \implies V_1 > V_\gamma$$
	che è una conclusione assurda. \\
	Allora il transistore $T1$ deve essere spento. \\
	Analogamente, anche $T2$ dovrà essere spento.
	Ma allora deve essere:
	$$\begin{cases} I_{C1} = 0 \\ I_{C2} = 0 \\ I_{RC} = I_{C1} + I_{C2} = 0 \\ V_u = V_{CC} - R_C I_{RC} \end{cases} \implies V_u = V_{CC}$$
	Quindi, quando entrambi i transistori sono spenti, il circuito si comporta come un invertitore.
	\subsubsection{$V_1 = V_L$  ;  $V_2 = V_H$}
	\label{subsub:gen-funz+ing-c1}
	Per quanto appena detto, il transistore $T1$ deve essere spento. \\
	Ma se $T1$ è spento ($I_{C1} = 0$), abbiamo un solo invertitore RTL ($T2$) che ha in ingresso un valore alto. \\
	Allora il transistore $T2$ è saturo. \\
	Per verificarlo, accertiamoci che sia:
	$$I_{C2} < \beta_F I_{B2}$$
	Per la Legge di Ohm è:
	$$I_{RC} = \frac{V_{CC} - V_{CESat}}{R_C}$$
	$$I_{RC} = I_{C1} + I_{C2} = I_{C2}$$
	$$I_{B2} = \frac{V_{CC} - V_\gamma}{R_B}$$
	Dunque, deve essere:
	$$\frac{V_{CC} - V_{CESat}}{R_C} < \beta_F \cdot \frac{V_{CC} - V_\gamma}{R_B}$$
	Questa condizione è verificata per:
	$$V_{CC} > V_\gamma + \frac{V_{CC} - V_{CESat}}{|A_V|}$$
	che è la condizione per avere un buon invertitore.\\\\
	Allora $V_u = V_L$.
	\subsubsection{$V_1 = V_H$  ;  $V_2 = V_L$}
	In questo caso abbiamo la condizione inversa a quella precedente. \\
	Allora avremo:
	$$V_u = V_L$$
	\subsubsection{$V_1 = V_H$  ;  $V_2 = V_H$}
	Siamo nella condizione in cui entrambi gli invertitori sono accesi. \\
	Verifichiamo che siano in regione di saturazione. \\
	In saturazione deve essere:
	$$\begin{cases} I_{C1} < \beta_F I_{B1} \\ I_{C2} < \beta_F I_{B2} \end{cases}$$
	$$I_{C1} + I_{C2} = \beta_F (I_{B1} + I_{B2})$$
	Per la Legge di Kirchhoff allora:
	$$\begin{cases} I_{RC} < \beta_F (I_{B1} + I_{B2}) \\ I_{RC} = \frac{V_{CC} - V_{CESat}}{R_C} \\ I_{B1} = \frac{V_{CC} - V_\gamma}{R_B} \\ I_{B2} = \frac{V_{CC} - V_\gamma}{R_B} \end{cases}$$
	$$\frac{V_{CC} - V_{CESat}}{R_C} < \frac{2\beta_F}{R_B} (V_{CC} - V_\gamma)$$
	che è certamente verificata per quanto detto nella sottosezione \ref{subsub:gen-funz+ing-c1}.\\\\
	Allora l’uscita deve essere:
	$$V_u = V_{CE1} = V_{CE2} = V_{CESat} = V_L$$
	\subsubsection{Considerazioni finali}
	Riportando la Tabella \ref{tab:table-9} in formato digitale notiamo che abbiamo realizzato una porta logica NOR a due ingressi.\\\\
	Possiamo allora generalizzare quanto appena ottenuto: volendo una porta NOR a tre ingressi con la semplice aggiunta di un terzo ramo in parallelo ai primi due. \\
	Inoltre la progettazione non è critica poiché se funziona il BJT, funziona anche la porta NOR.\\\\
	Ancora, avendo la funzione NOR, per la Legge di De Morgan, possiamo realizzare qualsiasi funzione logica di cui possiamo avere necessità.
	\begin{table}[ht]
		\centering
		\begin{tabular}{|c|c|c|}
			\hline
			$V_1$ & $V_2$ & $V_u$ \\
			\hline
			0 & 0 & 1 \\
			\hline
			0 & 1 & 0 \\
			\hline
			1 & 0 & 0 \\
			\hline
			1 & 1 & 0 \\
			\hline
		\end{tabular}
		\caption{Tabella di funzionamento del circuito di Figura \ref{fig:im-85} in formato digitale}
		\label{tab:table-9}
	\end{table}
	\subsection{Generalizzazione a qualsiasi invertitore}
	Assumiamo di avere un invertitore che presenti la caratteristica rossa di Figura \ref{fig:im-86} cioè una caratteristica in cui si abbiano due zone in cui il guadagno è inferiore ad uno ed una zona con guadagno maggiore di uno.\\\\
	Prima di tutto vogliamo imporre che il valore alto dell’ingresso sia pari al valore alto dell’uscita e che il valore basso in ingresso sia pari al valore alto dell’uscita.\\\\
	Consideriamo ora la serie di invertitori di Figura \ref{fig:im-87}.
	\begin{figure}[h]
		
		\begin{subfigure}{0.5\textwidth}
			\includegraphics[width=0.9\linewidth]{86.png} 
			\caption{Caratteristica di una invertitori in serie}
			\label{fig:im-86}
		\end{subfigure}
		\begin{subfigure}{0.5\textwidth}
			\includegraphics[width=0.9\linewidth]{87.png}
			\caption{Invertitori in serie}
			\label{fig:im-87}
		\end{subfigure}
		
		\caption{}
		\label{fig:image11}
	\end{figure}
	In una rete di questo tipo, gli ingressi degli invertitori pari coincidono con le uscite degli invertitori dispari e gli ingressi degli invertitori dispari coincidono con le uscite degli invertitori pari.\\\\
	Allora la caratteristica di entrambi è data dai punti di intersezione di quella di Figura \ref{fig:im-86}.\\\\
	Dall’intersezione della caratteristiche degli invertitori pari con la caratteristica degli invertitori dispari, ricaviamo i valori nominali e notiamo che abbiamo tre sole condizioni di funzionamento (indicate dai punti $A$, $B$ e $C$). \\
	I punti $A$ e $B$ sono inoltre simmetrici rispetto alla diagonale del primo e terzo quadrante e sono individuati da:
	$$\begin{cases} A = (V_L, V_H) \\ B = (V_H, V_L) \end{cases}$$
	Il punto $C$, che necessariamente è posto sulla diagonale, indica una terza condizione. \\
	Questo punto di lavoro individuato da:
	$$C = (V_{TL}, V_{TL})$$
	è la condizione in cui ciò che entra nell’invertitore è pari a ciò che esce dall’invertitore stesso.\\\\
	Supponiamo di porre in ingresso alla rete logica il valore $V_{TL}$. \\
	In assenza di rumore avremmo in uscita il valore $V_{TL}$. \\
	Tuttavia è impossibile mantenere costante il valore $V_{TL}$.\\\\
	Sfruttando il grafico della caratteristica notiamo facilmente che, anche a causa di un piccolo spostamento da $V_{TL}$, il valore dell’uscita si è avvicinato sempre più al valore nominale ($V_H$ per valori appena maggiori di $V_{TL}$ o $V_L$ per valori appena inferiori a $V_{TL}$). \\
	In altre parole, la qualità del segnale migliora dopo ogni inversione. \\
	Si parla in questo caso di \textbf{capacità di rigenerazione del segnale}.\\\\
	In particolare $V_{TL}$ è detto \textbf{valore di soglia logica} (Logic Threshold) poiché è il valore di soglia per cui il circuito è in grado di riportare un segnale rumoroso al suo valore nominale. \\
	Si noti che tale valore non è necessariamente posto a metà dell’escursione.\\\\
	Inoltre, poiché $V_{TL}$ è pressoché impossibile da mantenere, non dovremo preoccuparci della condizione di lavoro in $V_{TL}$.\\\\
	Si noti che la capacità di rigenerare il segnale dipende dalla pendenza elevata del tratto di curva in cui si posiziona $V_{TL}$. \\
	Infatti, per una pendenza inferiore ad uno, troviamo che in risposta ad una perturbazione di $V_{TL}$ il circuito riporta il segnale a $V_{TL}$. \\
	Ciò risulta evidente dal grafico di Figura \ref{fig:im-88} avendo un solo punto di intersezione esso deve, necessariamente, essere stabile.
	\begin{center}
		\begin{figure}[h]
			\centering
			\includegraphics[scale=0.45]{88.png}
			\caption{Se si ha una curva con guadagno $A_V < 1$ nella zona di massima pendenza si ottiene un circuito che tende a riportare l’uscita a $V_{TL}$}
			\label{fig:im-88}
		\end{figure}
	\end{center}
	\subsection{Generalizzazione di $N_M$}
	Generalizziamo ora il concetto di margine di immunità ai disturbi. Riprendiamo la caratteristica di Figura \ref{fig:im-89}.
	\begin{center}
		\begin{figure}[h]
			\centering
			\includegraphics[scale=0.45]{89.png}
			\caption{Caratteristica di un invertitore. \\ Sono evidenziati i margini di tolleranza al rumore.}
			\label{fig:im-89}
		\end{figure}
	\end{center}
	In questo caso la determinazione di $N_M$ è più complessa rispetto a quanto visto in precedenza poiché non abbiamo tratti a pendenza nulla. \\
	Possiamo tuttavia riconoscere due tratti a pendenza:
	$$|A_V| < 1$$
	in cui il rumore viene attenuato ed un tratto a pendenza:
	$$|A_V| > 1$$
	in cui il rumore viene amplificato.\\\\
	Possiamo etichettare i due punto in cui il guadagno è $A_V = -1$ come:
	$$(V_{iL \ Max}, V_{OH \ Min})$$
	per l’ingresso di un segnale basso. \\
	Si noti che $V_{OH \ Min}$ è più vicino a $V_H$ di quanto non lo sia $V_{iL \ Min}$ a $V_L$:
	$$V_H - V_{OH \ Min} > V_L - V_{iL \ Max}$$
	Per quanto riguarda l’ingresso alto abbiamo invece:
	$$(V_{iH \ Min}, V_{OL \ Max})$$
	$$V_H - V_{iH \ Min} < V_L - V_{OL \ Max}$$
	Risulta inoltre:
	$$\begin{cases} V_i < V_{iL \, Max} \Longrightarrow V_u > V_{OH \, Min} \\ V_i > V_{iH \, Min} \Longrightarrow V_u < V_{OL \, Max} \end{cases}$$
	Dal grafico di Figura \ref{fig:im-90} notiamo che il margine di immunità ai disturbi associato al valore alto è:
	$$N_{MH} = V_{OH \, Min} - V_{iH \, Min}$$
	e, analogamente, per il margine basso è:
	$$N_{ML} = V_{iL \, Max} - V_{OL \, Max}$$
	\begin{center}
		\begin{figure}[h]
			\centering
			\includegraphics[scale=0.45]{90.png}
			\caption{Distribuzione dei margini di immunità ai disturbi}
			\label{fig:im-90}
		\end{figure}
	\end{center}
	Allora il margine di immunità ai disturbi complessivo è:
	$$N_M = \min \left\{N_{ML}, N_{MH}\right\}$$
	Ancora, vogliamo che entrambi i margini siano positivi. \\
	Dunque è:
	$$V_{OH\ Min} - V_{iH\ Min} + V_{iL\ Max} - V_{OL\ Max} > 0$$
	da cui segue che:
	$$\underbrace{V_{OH\ Min} - V_{OL\ Max}}_{\Delta V_u} > \underbrace{V_{iH\ Min} - V_{iL\ Max}}_{\Delta V_i}$$
	In questo caso non possiamo approssimare il rapporto incrementale alla derivata. \\
	Tuttavia, risulta comunque evidente che deve essere:
	$$\frac{\Delta V_u}{\Delta V_i} > 1$$
	che conferma quanto detto all’inizio dell’analisi.\\\\
	Anche in questo caso generico è evidente che la somma $N_{ML} + N_{MH}$ è un valore finito. \\
	Inoltre è possibile individuare graficamente $N_{MH}$ e $N_{ML}$ utilizzando un compasso.
	\subsection{Incongruenza}
	Tutti i ragionamenti fatti fino a questo punto si basano sull’ipotesi che la rete sia costituita dalla serie di invertitori di Figura \ref{fig:im-87}.
	Per l’analisi abbiamo però utilizzato le caratteristiche degli invertitori calcolate a vuoto. \\
	Quando costruiamo il circuito di Figura \ref{fig:im-87}, tuttavia, connettiamo in serie più invertitori. \\
	Dunque, per il primo invertitore è:
	$$V_u = V_{CC} - R_C \left(I_C - I’_B\right)$$
	dove $I’_B$ è la corrente di base dovuta al secondo invertitore.\\\\
	Analizzeremo la caratteristica, considerando correttamente la condizione di carico, nelle prossime sezioni.
	\subsection{Rete di invertitori}
	Consideriamo anzitutto una rete composta da due invertitori RTL identici come in Figura \ref{fig:im-91}.
	\begin{center}
		\begin{figure}[h]
			\centering
			\includegraphics[scale=0.45]{91.png}
			\caption{Due invertitori RTL in serie}
			\label{fig:im-91}
		\end{figure}
	\end{center}
	In particolare la corrente sul nodo evidenziato (il nodo di connessione tra le due porte) è detta \textbf{corrente di fan out} poiché il quel nodo la corrente si apre a ventaglio.\\\\
	Nel caso di fan out pari a zero la caratteristica è quella riportata in Figura \ref{fig:im-92}.
	\begin{center}
		\begin{figure}[h]
			\centering
			\includegraphics[scale=0.45]{92.png}
			\caption{Caratteristica di un invertitore RTL a vuoto}
			\label{fig:im-92}
		\end{figure}
	\end{center}
	Dal momento che abbiamo la connessione di due transistori dovremmo valutare nove condizioni. \\
	Tuttavia poiché il secondo BJT dipende dal primo possiamo ridurre il numero di casi di studio.
\end{document}